% =====================================================================
%  BIBLIOGRAFÍA (versión mejorada, APA 7: autores corporativos, fechas, sin
%  «Obtenido de»; entradas citadas pero ausentes añadidas; entradas sin cita
%  retiradas al final como comentario; dudas marcadas con «VERIFICAR»).
%  Cada referencia es un \item dentro del entorno referenciasapa.
% =====================================================================
\chapter*{BIBLIOGRAFÍA}
\addcontentsline{toc}{chapter}{BIBLIOGRAFÍA}
\label{cap:bibliografia}

\begin{referenciasapa}

\item Abrahamsson, P., Hanhineva, A., Hulkko, H., Ihme, T., Jäälinoja, J., Korkala, M., Koskela, J., Kyllönen, P., \& Salo, O. (2004). Mobile-D: An agile approach for mobile application development. En J. Vlissides (Ed.), \textit{OOPSLA'04: Companion to the 19th annual ACM SIGPLAN conference on object-oriented programming, systems, languages, and applications} (pp. 174--175). ACM. \url{https://doi.org/10.1145/1028664.1028736}

\item Abstracta Team. (2024). \textit{Metodologías de desarrollo de software: guía para optimizar tu proyecto}. Abstracta. \url{https://es.abstracta.us/blog/metodologias-de-desarrollo-de-software-guia-completa/}

\item Amazon Web Services. (2025). \textit{¿Qué es la IA? Explicación de la inteligencia artificial}. \url{https://aws.amazon.com/es/what-is/artificial-intelligence/}

\item Android Developers. (2024). \textit{Introducción a Android Studio}. \url{https://developer.android.com/studio/intro?hl=es-419}

\item Apple Inc. (2025). \textit{ARKit 6: Augmented reality}. Apple Developer. \url{https://developer.apple.com/augmented-reality/arkit/}
% VERIFICAR URL

\item Apple. (2023). \textit{Apple Developer Documentation}. \url{https://developer.apple.com/}

\item AppMaster. (2023, 4 de septiembre). \textit{Programación móvil}. \url{https://appmaster.io/es/glossary/programacion-movil}

\item Apps, L. B. (2011). \textit{Guía de aplicaciones móviles}. Fundación Telefónica.
% VERIFICAR autor y editorial

\item Atlassian. (2023). \textit{About us: Trello history, logos \& customers | Trello}. \url{https://trello.com/about}

\item Azuma, R. T. (1997). A survey of augmented reality. \textit{Presence: Teleoperators and Virtual Environments, 6}(4), 355--385. \url{https://doi.org/10.1162/pres.1997.6.4.355}

\item B, G. (2025, 13 de marzo). \textit{Qué es GitHub y cómo empezar a usarlo}. ES Tutoriales. \url{https://www.hostinger.com/es/tutoriales/que-es-github}
% VERIFICAR nombre completo del autor

\item Beck, K., \& Andres, C. (2005). Extreme Programming Explained: Embrace change. En \textit{Extreme Programming Explained.} Addison-Wesley.

\item Cruz. (2024). \pendiente{completar: apellido y nombre del autor, título del proyecto de grado, Universidad Mayor de San Simón, ciudad}

\item Cuadrado, G. C. (2022, 22 de julio). \textit{Qué es Visual Studio Code y qué ventajas ofrece}. OpenWebinars. \url{https://openwebinars.net/blog/que-es-visual-studio-code-y-que-ventajas-ofrece/}

\item Deloitte. (2019). \textit{¿Qué es React Native?} Deloitte España. \url{https://www.deloitte.com/es/es/services/consulting/blogs/todo-tecnologia/que-es-react-native.html}

\item DuBois, P. (2013). \textit{MySQL.} Addison-Wesley Professional.

\item Fernández. (2018). \pendiente{completar: autor, título del proyecto de grado, universidad, ciudad; unificar la grafía «Fernández» en el texto (aparece también «Fernandez»)}

\item Flutter. (s. f.). \textit{Flutter: crea aplicaciones para cualquier pantalla}. \url{https://esflutter.dev/}

\item Gabriel Maide, E., \& Pacienzia, J. (2020, noviembre). \textit{Metodologías de desarrollo de software aplicadas a proyectos de automatización de procesos}. Pistas Educativas. \url{https://pistaseducativas.celaya.tecnm.mx/index.php/pistas/article/viewFile/2381/1900}

\item Gironés, T. (2012). \textit{Aplicaciones móviles: Conceptos y desarrollo.}

\item Git. (s. f.). \textit{Git}. \url{http://git-scm.com/}

\item Gjestvang, C., Abrahamsen, F., Stensrud, T., \& Haakstad, L. A. H. (2020). Motives and barriers to initiation and sustained exercise adherence in a fitness club setting---A one-year follow-up study. \textit{Scandinavian Journal of Medicine \& Science in Sports, 30}(9), 1796--1805. \url{https://doi.org/10.1111/sms.13736}
% CONFIRMAR DOI; la referencia original mezclaba este artículo con otro de Frontiers in Psychology y un enlace de PsycNET

\item Gjestvang, C., Stensrud, T., Paulsen, G., \& Haakstad, L. A. H. (2021). Stay true to your workout: Does repeated physical testing boost exercise attendance? A one-year follow-up study. \textit{Journal of Sports Science and Medicine, 20}, 35--44.

\item Gómez, A. A. (2015). \textit{Comparación de dos tecnologías de desarrollo de aplicaciones móviles desde la perspectiva de los atributos de calidad} [Tesis de maestría]. Universidad Nacional de Colombia.
% VERIFICAR tipo de trabajo

\item Google Cloud. (2025). \textit{Inteligencia artificial (IA): una guía fácil de entender}. \url{https://cloud.google.com/learn/what-is-artificial-intelligence?hl=es-419}

\item Google for Developers. (2025). \textit{ARCore: Crea nuevas experiencias de realidad aumentada}. \url{https://developers.google.com/ar?hl=es-419}

\item Google. (2023). \textit{Android Developers -- Official Documentation}. \url{https://developer.android.com/}

\item Group, L. F. (2021). \textit{Ejercicios en el gimnasio}. Google Play. \url{https://play.google.com/store/apps/details?id=gymworkout.gym.gymlog.gymtrainer&hl=es_BO}
% VERIFICAR autor corporativo real y ajustar la cita (Ejercicios en el Gimnasio, 2021)

\item Gurnov, A., \& Wrike. (2025, 13 de noviembre). \textit{Las 21 mejores herramientas de gestión de proyectos}. Wrike. \url{https://www.wrike.com/es/project-management-guide/faq/que-son-las-herramientas-de-gestion-de-proyectos/}

\item Hayes, M., \& Downie, A. (2025, 28 de noviembre). \textit{Realidad aumentada}. IBM. \url{https://www.ibm.com/es-es/think/topics/augmented-reality}

\item Huang, S. (2020). \textit{NoSQL Databases: Concepts and Applications.} Cambridge: Morgan Kaufmann.

\item IBM. (2025). \textit{¿Qué son las redes neuronales?} \url{https://www.ibm.com/es-es/think/topics/neural-networks}

\item IEEE. (1990). \textit{IEEE standard glossary of software engineering terminology} (IEEE Std 610.12-1990). Institute of Electrical and Electronics Engineers.

\item Ionic. (2023). \textit{Ionic Framework}. Ionic Documentation. \url{https://ionicframework.com/docs}

\item Isaac. (2025). \textit{Firebase: qué es y cómo funciona la plataforma de Google}. Mundobytes. \url{https://mundobytes.com/firebase-que-es/}

\item Kolt, G., \& Snyder-Mackler, L. (2007). \textit{Terapias físicas en el deporte y el ejercicio}.
% VERIFICAR edición, traductor y editorial; no se abrió la fuente

\item Kraemer, W. J., \& Ratamess, N. A. (2004). Fundamentals of resistance training: Progression and exercise prescription. \textit{Medicine \& Science in Sports \& Exercise, 36}(4), 674--688. \url{https://doi.org/10.1249/01.MSS.0000121945.36635.61}

\item Laux, E., \& S., J. (2016). \textit{BodBot AI entrenador personal}. \url{https://www.bodbot.com/}
% VERIFICAR autores y ajustar la cita (BodBot IA Entrenador Personal, 2016)

\item Leyva, R., León, G., \& García, F. (2016). Metodología Mobile-D para el desarrollo de aplicaciones móviles. \textit{Revista Cubana de Ciencias Informáticas, 10}(3), 45--57.
% VERIFICAR: la referencia original enlazaba una figura de ResearchGate, no el artículo

\item López Mora, S. (2025, noviembre). \textit{Firebase: qué es, para qué sirve, funcionalidades y ventajas}. DIGITAL55. \url{https://digital55.com/blog/que-es-firebase-funcionalidades-ventajas-conclusiones/}

\item Meier, R., \& Lake, I. (2018). \textit{Professional Android} (4.ª ed.).
% VERIFICAR editorial

\item Microsoft. (2022). \textit{Xamarin documentation}. Microsoft Learn. \url{https://learn.microsoft.com/en-us/xamarin/}

\item Mifflin, M. D., St Jeor, S. T., Hill, L. A., Scott, B. J., Daugherty, S. A., \& Koh, Y. O. (1990). A new predictive equation for resting energy expenditure in healthy individuals. \textit{The American Journal of Clinical Nutrition, 51}(2), 241--247. \url{https://doi.org/10.1093/ajcn/51.2.241}
% CONFIRMAR DOI

\item Mitchell, T. (1997). \textit{Machine Learning.} New York: McGraw-Hill.

\item MongoDB, Inc. (2024). \textit{MongoDB Atlas documentation}. \url{https://www.mongodb.com/docs/atlas/}

\item Navarro, F., Martínez, A., \& Martínez, J. (2018). \textit{Realidad Virtual y Realidad Aumentada.} Ra-Ma.

\item Nielsen, J. (1993). \textit{Usability Engineering.} Academic Press.

\item Open Handset Alliance. (2022). \textit{Android Open Source Project overview}. \url{https://source.android.com/?hl=es-419}

\item Oracle Corporation. (2024). \textit{MySQL documentation}. MySQL. \url{https://dev.mysql.com/doc/}

\item Organización Mundial de la Salud. (2000). \textit{Obesity: Preventing and managing the global epidemic. Report of a WHO consultation} (WHO Technical Report Series, 894). Ginebra: Organización Mundial de la Salud.

\item Orientanet. (s. f.). \textit{¿Qué es entrenamiento físico y tipos?} \url{https://www.orientanet.es/que-es-entrenamiento-fisico-y-tipos/}

\item Pressman, R. S. (2010). \textit{Software Engineering: A Practitioner's Approach.} Higher Education.

\item Pykes, K. (2024). \textit{Aprendizaje automático en Python: tutorial de Scikit-Learn}. DataCamp. \url{https://www.datacamp.com/es/tutorial/machine-learning-python}

\item React Native. (2025). \textit{Core components and native components}. Meta. \url{https://reactnative.dev/docs/intro-react-native-components}
% VERIFICAR título exacto y ajustar la cita (React Native, 2025)

\item Sandler, M., Howard, A., Zhu, M., Zhmoginov, A., \& Chen, L.-C. (2018). MobileNetV2: Inverted residuals and linear bottlenecks. \textit{Proceedings of the IEEE Conference on Computer Vision and Pattern Recognition}, 4510--4520. \url{https://openaccess.thecvf.com/content_cvpr_2018/html/Sandler_MobileNetV2_Inverted_Residuals_CVPR_2018_paper.html}

\item Santander Open Academy. (2020, 21 de diciembre). \textit{Metodologías de desarrollo de software: ¿qué son?} \url{https://www.santanderopenacademy.com/es/blog/metodologias-desarrollo-software.html}

\item Sanz, D. (2019). \textit{Introducción a Android}. \url{https://dspace.itsjapon.edu.ec/jspui/bitstream/123456789/434/1/introduccion-android.pdf}

\item Schoenfeld, B. J. (2010). The mechanisms of muscle hypertrophy and their application to resistance training. \textit{Journal of Strength and Conditioning Research, 24}(10), 2857--2872. \url{https://doi.org/10.1519/JSC.0b013e3181e840f3}

\item Schwaber, K., \& Sutherland, J. (2020). \textit{The Scrum Guide: The Definitive Guide to Scrum.} Boston: Scrum.org.

\item Silva Amado, H. S. (2024). \textit{Sistema de visión artificial para la identificación de ejecución de ejercicios físicos} [Trabajo de grado, Universidad Santo Tomás]. Tunja, Colombia.

\item SQLite Consortium. (2024). \textit{SQLite home page}. SQLite.org. \url{https://sqlite.org/}

\item StatCounter Global Stats. (2025, noviembre). \textit{Mobile operating system market share Plurinational State of Bolivia}. \url{https://gs.statcounter.com/os-market-share/mobile/bolivia}

\item Tomás, J. (2019). \textit{El gran libro de Android} (Vol. 7). Bogotá: Alfaomega.

\item U.S. Department of Agriculture. (s. f.). \textit{FoodData Central}. \url{https://fdc.nal.usda.gov/}

\item Viera, U. (s. f.). \textit{Introducción a React Native: crea apps móviles nativas en iOS y Android}. urianviera.com. \url{https://urianviera.com/react-native/introduccion-a-react-native}

\item wger Project. (s. f.). \textit{wger Workout Manager: REST API}. \url{https://wger.de/en/software/api}

\item WorkoutX. (s. f.). \textit{WorkoutX API: exercise database and GIF animations}. \url{https://workoutxapp.com/}

\item Yépez, M. D. (2017). \textit{Club Ensayos}. \url{https://www.clubensayos.com/Negocios/Proyecto-de-grado-Creaci\%C3\%B3n-de-un-centro-de/4574087.html}

\item Yuhonas. (s. f.). \textit{free-exercise-db: open public domain exercise dataset}. GitHub. \url{https://github.com/yuhonas/free-exercise-db}

% ---------------------------------------------------------------------
% ENTRADAS SIN CITA EN EL TEXTO (retiradas de la lista; restaurar solo si se citan)
% ---------------------------------------------------------------------
% Booch, G., Maksimchuk, R., Engel, M., Young, B., Conallen, J., \& Houston, K. (2007). \textit{Object-Oriented Analysis and Design with Applications.} Addison-Wesley.
% \textit{Curso para desarrollo de apps con React Native}. (s.f.). Obtenido de Código Facilito: \url{https://codigofacilito.com/articulos/que-es-react-native}
% GeeksforGeeks. (11 de julio de 2025). \textit{Incremental Process Model Software Engineering.} Obtenido de GeeksforGeeks: \url{https://www.geeksforgeeks.org/software-engineering/software-engineering-incremental-process-model/}
% Gómez, \& Hernández. (2016). Mobile-D: Metodología ágil para el desarrollo de aplicaciones móviles. \textit{Revista Orbis}.
% Gómez, A. A. (2015). \textit{Comparación de dos tecnologías de desarrollo de aplicaciones móviles desde la perspectiva de los atributos de calidad}. Obtenido de ResearchGate: \url{https://www.researchgate.net/figure/Figura-5-Arquitectura-de-Android_fig3_319975215}
% Guide, I. S. (12 de septiembre de 2023). \textit{Modelo en espiral: modelo de reducción de riesgos para proyectos de software.} Obtenido de \url{https://www.ionos.mx/startupguide/productividad/modelo-en-espiral/}
% Nystrom, K. (2021). \textit{Kotlin Programming: The Big Nerd Ranch Guide. Big Nerd Ranch.} Pearson Education Store.
% Sharma, M. (2018). \textit{Productivity gains with modern IDEs: A review of IntelliJ-based environments.}
% Smith, J., \& Darrow, M. (2019). \textit{Gradle for Android Development.} O'Reilly Media.
% Sommerville, I. (2011). \textit{Software Engineering.} Boston: Addison-Wesley.

\end{referenciasapa}
